\section{The benchmark}\label{sec:benchmark}

% Describes the battery as specified in docs/scenarios.md at the commit of the battery run,
% scripts/bench.sh (run matrix, time limits) and tools/stats.py (statistics). Re-check against
% those files if the battery changes. Table 2 is transcribed from docs/scenarios.md §4: design
% parameters, not results.

The battery is a set of scenarios, each with a set-up, a reference solution and an error
measure, written once against an engine-independent interface and run on every engine with every
contact method.

\subsection{What a reference rests on: classes A, B and C}\label{sec:classes}

Every scenario is classified by what its reference rests on.
\begin{description}
\item[Class A] A rigid-body result that does not depend on the contact model: the acceleration
  of a block on an incline, the time at which a sliding ball starts to roll, the period of a
  ball rolling in a bowl, a wedge that must hold if and only if $\tan\alpha < \mu$. Only these
  scenarios compare a compliant law with a contact model of another kind on equal terms.
\item[Class B] A reference that assumes a pressure distribution: the stop time of a spinning
  disk, ring or cone, the coupled sliding and spinning of a disk, a journal in a bearing without
  clearance. The nominal reference uses the distribution that a rigid body produces on a Winkler
  foundation (uniform for flat faces, a cosine for the journal), so agreement with it is in part
  agreement with the model under test. Other distributions are defensible: uniform wear, the
  rigid punch on an elastic half-space, a line contact. The part of the reference that does not
  depend on the model is the \emph{band} spanned by those distributions: for the disk, the
  friction torque over $\mu N R$ is $2/3$ for uniform pressure, $1/2$ for uniform wear and
  $\pi/4$ for the rigid punch. A method inside the band is not wrong.
\item[Class C] Consistency with the Winkler foundation. The reference contains $K$ or
  integrates the law's own force: the penetrations of a stack, the tilt of an eccentric load,
  the angle at which a tall box tips over a soft support, the rebound of an impact governed by
  the law's own differential equation, the depth reached in a press fit. These scenarios say
  whether a method is a Winkler foundation of stiffness $K$. They do not say whether a method is
  accurate, and they are never used to compare the laws with contact models of another kind.
\end{description}

\subsection{The scenarios}\label{sec:scenarios}

Table~\ref{tab:scenarios} lists the \result{battery/scenarios} scenarios with the quantity that
each one measures; Appendix~\ref{app:scenarios} gives their parameters and references, and the
specification in the repository the dimensions, the mesh criteria and the derivation of each
reference. Four rules hold for all of them. Parameters are fixed before any run, and a scenario
is not changed because a method scores badly on it; the foundations have $K$ between $10^6$ and
$10^8$\,N/m$^3$ and $h = 5$\,s/m. Initial conditions belong to no model: a body that starts in
contact penetrates 1\,\textmu m whatever the method. The reference covers only what it can: the
scenarios of a body rolling on a plane measure what the rigid theory fixes whatever the rolling
resistance of the foundation is (the acceleration at zero speed, the onset of rolling). And the
curved meshes approximate the solid from the inside, with 64 segments per turn in the base
variant; each flat cap of a cylinder, a disk or a frustum is one polygon triangulated as a fan
from one of its rim vertices, and each planar ground is a slab with two triangles per face.

\begin{table}[p]
\caption{The scenarios, by class. $^{2}$: second set (Section~\ref{sec:variants}). For class B
the band of admissible pressure distributions is given as the factor on the nominal stop time}
\label{tab:scenarios}
\footnotesize
\begin{tabular}{@{}lp{0.47\linewidth}p{0.36\linewidth}@{}}
\toprule
Scenario & Set-up & Measured against the reference \\
\midrule
\multicolumn{3}{@{}l}{\emph{Class A: rigid-body references}} \\
\scn{block} & Box sliding down a $30^\circ$ incline & Final velocity \\
\scn{cylinder} & Cylinder rolling down $15^\circ$ from rest & Acceleration at zero speed, slip \\
\scn{sphere} & Sphere rolling down $20^\circ$ from rest & Acceleration at zero speed, slip \\
\scn{push} & Box pushed by a slowly growing force & Onset of sliding, velocity, creep before it \\
\scn{billiard} & Ball launched sliding, without spin & Time and velocity at which rolling starts \\
\scn{vgroove} & Box sliding in a $90^\circ$ groove tilted $30^\circ$ & Final velocity \\
\scn{twoblocks} & Block on a block on a $30^\circ$ incline & Common velocity, no relative slip \\
\scn{bowl} & Ball rolling in a spherical bowl & Period, amplitude (constant) \\
\scn{trough} & Cylinder rolling in a cylindrical trough & Period, amplitude (constant) \\
\scn{journal} & Shaft spinning in a hole with clearance & Stop time, angle of climb \\
\scn{pieces1}, \scn{2}, \scn{8}, \scn{64}$^{2}$ & Sliding and spinning boxes on a table cut into 1, 2, 8 and 64 pieces & Travel, stop time, no drift; same final poses on every cut \\
\scn{wedge}$^{2}$ & Two blocks on a plane tilted slowly; one must hold, one slide & Creep, residual speed, sliding velocity \\
\scn{oblique}$^{2}$ & Spheres hitting a plane obliquely, sliding and sticking & Velocities and spins after the impact \\
\scn{edge}$^{2}$ & Box sliding off the edge of a table & Times at which it has pitched 0.1 and 0.25\,rad \\
\scn{turntable}$^{2}$ & Ball rolling on a free rotating table & Orbit rate, orbit radius (constant) \\
\scn{energy}$^{2}$ & Frictionless, undamped ball sliding in a faceted bowl & Period, drift of the energy \\
\scn{tower}$^{2}$ & Ten stacked blocks, gravity tilted to $5^\circ$ & Nothing moves: creep, residual speed \\
\scn{massratio}$^{2}$ & A 1\,kg block under a 1000\,kg block & Nothing moves: creep, residual speed \\
\midrule
\multicolumn{3}{@{}l}{\emph{Class B: references that assume a pressure distribution}} \\
\scn{disk} & Disk spinning flat on a plane & Stop time (band 0.85--1.33) \\
\scn{contensou} & Disk sliding and spinning & Stop time (band 0.95--1.12) \\
\scn{ring} & Annulus spinning flat & Stop time (band 1--1.02) \\
\scn{bearing} & Shaft spinning in a hole of its own radius & Stop time (band 0.81--1.27) \\
\scn{cone} & Frustum spinning in a conical seat & Stop time (band 1--1.04) \\
\midrule
\multicolumn{3}{@{}l}{\emph{Class C: consistency with the Winkler foundation}} \\
\scn{tip} & Tall box, gravity tilted slowly & Tipping angle on the foundation \\
\scn{drop} & Box dropped flat & Rebound velocity of the law's impact equation \\
\scn{hole} & Cube entering a square hole with interference & Depth at which friction stops it \\
\scn{corner} & Cube pressed into a corner & Penetrations on the three faces \\
\scn{offset} & Plate with an eccentric centre of mass & Tilt at rest \\
\scn{stack} & Three cubes, gravity tilted to $5^\circ$ & Compressions; creep \\
\scn{restitution}$^{2}$ & Spheres hitting a plane at four speeds & Rebounds of the law's impact equation \\
\botrule
\end{tabular}
\end{table}

\subsection{The error of a run}\label{sec:error}

The error of a run is the worst normalized deviation of every quantity that the reference
fixes, and not of the main result alone, so that a box that creeps before the onset of sliding
and then reaches the right velocity is not scored as correct. A value with a non-zero reference
is normalized by that reference, and a quantity that the reference keeps at zero (creep, drift,
relative slip, residual speed) by a length or a speed of the body. A run \emph{holds} when it
finishes with an error below 100\,\%; it is \emph{lost} when the error is larger (bodies thrown
apart or sliding away), and it has \emph{no value} when it diverges, exceeds its time limit or
never produces a quantity of the reference (a body that never stopped or never tipped). The
battery records the size of each term and not its sign, so for class B a run is known to be
inside the band when its error is below the narrow side of the band, outside when it is above
the wide side, and undetermined in between. Errors are not commensurable across scenarios:
methods are compared within a scenario, and summaries across scenarios use counts at fixed
thresholds and ratios of the errors of two methods on the same scenario, never means of errors.

\subsection{Variants, methods and native models}\label{sec:variants}\label{sec:natives}

Every scenario is run in a base variant (64 segments per turn, $\dt = 1$\,ms) and in series:
the step ($\dt = 2$, 0.5, 0.2 and 0.1\,ms), the mesh (32 and 128 segments per turn, for the
scenarios with curved surfaces), the stiffness ($K\times10$ at $\dt = 0.2$\,ms, not for class
C) and eight perturbations of the initial velocities. The matrix is run with the six methods of
the library on each of the three engines (\law{gonthier}, \law{gonthier-thesis},
\law{distributed}, its two friction quadratures and \law{pcm}) and with six native models, each
with a wall-clock limit. Twenty-one scenarios were written while the laws were being developed;
the other twelve, the \emph{second set}, were written with their references before any law was
run on them, and the laws were changed after their first run (Section~\ref{sec:limitations}).

The native models stand for three families, and the results are read by family. \emph{Point
contact}: Drake's compliant point contact, Chrono's penalty model on a collision manifold (SMC)
and MuJoCo's soft constraints with its recommended and with its default parameters; a contact
between two faces is a few points with a stiffness in N/m, so a pressure distribution over the
face is not part of the model. \emph{Complementarity contact}: Chrono NSC, rigid contact with
Coulomb friction, in which $K$ and $h$ do not enter, so class C does not apply to it.
\emph{Pressure-field contact}: Drake hydroelastic, configured with a modulus $E = K H$ so that
the pressure is $K$ times the depth; it is then the same pressure model as the two laws near
the surface, evaluated on tetrahedral fields, and its comparison with \law{distributed} is one
of two discretizations. Each model has one configuration for every scenario, built from the
engine's documentation \cite{DrakeHydroGuide,DrakeHydroParams,DrakePointContact,DrakeFoamBrick,
ChronoCollisions,ChronoSMCSource,ChronoSystem,Barzilai1988,MuJoCoDocs,MuJoCoModeling}; the
configurations are ours and were not reviewed by the developers of the engines, so the main
text draws from the native runs only what follows from a property of the model family, and
Section~S3 gives the configurations and every other result, labelled as results of our
configurations. In the tables the methods are abbreviated Gon (\law{gonthier}), G18
(\law{gonthier-thesis}), Dis (\law{distributed}), Dis-v and Dis-1 (its vertex and one-point
friction rules), PCM, Dr-h and Dr-p (Drake hydroelastic and point), Ch-S and Ch-N (Chrono SMC
and NSC) and Mj-s and Mj-d (MuJoCo soft and default); the methods of the library are shown by
their runs on Chrono.
